\documentclass[12pt]{article}
\usepackage[utf8]{inputenc}
\usepackage[T1]{fontenc}
\usepackage{hyperref}
\title{Tanaina Language Resources for K-12 Curriculum Development}
\author{Compiled by Claude Sonnet 5 (claude-sonnet-5) for Andr\'as Kornai}
\date{}
\begin{document}
\maketitle
\section{Introduction}
Tanaina, better known as Dena'ina Qenaga (ISO 639-3 tfn), is an Athabaskan language of south-central Alaska. Speaker counts in the sources differ. Estimates range from about 50 first-language speakers to fewer than 100 fluent speakers. A Kenaitze page reports only five to seven first-language speakers alive when its dictionary was published.
\section{Primary Educational Resources}
\begin{itemize}
\item Albert Wassillie, \emph{Dena'ina Athabaskan Junior Dictionary} (1979, edited by James Kari, Alaska Native Language Center), an English-to-Dena'ina dictionary: \url{https://wiki.ubc.ca/Documentation:RelLex/Dena%27ina_Athabaskan_Junior_Dictionary:_Dena%27ina_Qenaga_Duch%27duldih}
\item Dena'ina Language Institute of the Kenaitze Indian Tribe, with learner resources: \url{https://www.kenaitze.org/education/denaina-language-institute/}
\end{itemize}
\section{Linguistic and Reference Materials}
\begin{itemize}
\item James Kari, \emph{Kenai Tanaina Noun Dictionary} (1974, preliminary, Outer Inlet dialect): \url{https://wiki.ubc.ca/Documentation:RelLex/Kenai_Tanaina_Noun_Dictionary}
\item \emph{Dena'ina Topical Dictionary} (Alaska Native Language Center), with excerpts online: \url{https://qenaga.org/learning/dtd}
\item Omniglot overview of the writing system: \url{https://www.omniglot.com/writing/denaina.htm}
\end{itemize}
\section{Religious and Community Resources}
\begin{itemize}
\item Kenaitze Indian Tribe, Dena'ina Language Institute: \url{https://www.kenaitze.org/education/denaina-language-institute/}
\item Alaska Native Languages page for Dena'ina: \url{https://www.alaskanativelanguages.org/dena-ina-qenaga-dena-ina}
\end{itemize}
\section{Digital and Audio Resources}
\begin{itemize}
\item Kenaitze audio dictionary and pocket dictionary (described in the Kenaitze page above).
\item Dena'ina site at Kenai Peninsula College: \url{https://denaina.kpc.alaska.edu}
\end{itemize}
\section{Curriculum Development Context}
Classes are reported at Kenai Peninsula College and the University of Alaska Anchorage. The search did not show a K-12 programme. Dialect differences (Inlet and Outer Inlet, among others) matter when choosing materials.
\end{document}